\chapter{Complexity — Seam}
%%% generated by the blueprint pipeline from TCSlib.Complexity.TuringMachine.Build.Contract.Seam
\section{Overview}

Declarations of the Lean module \texttt{TCSlib.Complexity.TuringMachine.Build.Contract.Seam}.

\section{Declarations}

\begin{lemma}[State after relabeling]
\lean{Turing.Contract.mapState_state}
\label{Turing.Contract.mapState_state}
\leanok
\uses{Turing.Cfg, Turing.Cfg.mapState}
\difficulty{1}
For any configuration $c$ and function $f$ between state sets, the state of the
configuration obtained by relabeling $c$ by $f$ is the image of $c$’s optional state
under $f$.
\end{lemma}

\begin{lemma}[State recovery after constant renaming]
\lean{Turing.Contract.mapState_const_inl}
\label{Turing.Contract.mapState_const_inl}
\leanok
\uses{Turing.Cfg, Turing.Cfg.mapState}
\difficulty{3}
Let $S_1,S_2$ be state types, let $e\in S_1$, and let $d$ be a configuration whose state
is the left injection of $e$ into the disjoint union $S_1\sqcup S_2$. First map every
state to $e$, then map $e$ back through the left injection; the resulting configuration
is exactly $d$.
\end{lemma}

\begin{lemma}[Trajectory of Sequential Machine Composition]
\lean{Turing.Contract.seq_traj}
\label{Turing.Contract.seq_traj}
\leanok
\uses{Turing.Cfg, Turing.Cfg.mapState, Turing.Contract, Turing.Contract.mapState_const_inl, Turing.Contract.mapState_state, Turing.MultiTapeTM, Turing.MultiTapeTM.runFrom, Turing.MultiTapeTM.runFrom_state_ne_none_of_le, Turing.seamCompTM, Turing.seamCompTM_runFrom_left, Turing.seamCompTM_run_ofCfg}
\difficulty{4}
Let $M_1$ and $M_2$ be multi-tape Turing machines with the same number of work tapes and
symbol alphabet, with state sets $S_1$ and $S_2$. Fix an entry state $e_1$ and an exit
state $f_1$ of $M_1$, an entry state $e_2$ of $M_2$, and a contract $C_1$ for $M_1$
whose entry is $e_1$ and whose exit set is exactly $\{\mathrm{some}(f_1)\}$. For a parameter $i$
and a configuration $d$ of the machine with state set $S_1\oplus S_2$, suppose $d$ is in
state $(e_1,\mathrm{inl})$ and replacing its state by $e_1$ gives a configuration
satisfying $C_1$'s precondition for $i$. At every time $u$ up to and including the
contract running time, the composed machine's run from $d$ equals the $M_1$ run from
that state-replaced configuration with states mapped into the left summand, and the
latter run has a nonempty state. For every $v$, the composed run at the contract running
time plus one plus $v$ equals the $M_2$ run from the contract endpoint with its state
replaced by $e_2$, mapped into the right summand; that state-replaced endpoint is in
state $e_2$.
\end{lemma}

\begin{definition}[Sequential composition of contracts]
\lean{Turing.Contract.seq}
\label{Turing.Contract.seq}
\leanok
\uses{Turing.Cfg.mapState, Turing.Contract, Turing.Contract.mapState_const_inl, Turing.Contract.mapState_state, Turing.Contract.seq_traj, Turing.MultiTapeTM, Turing.MultiTapeTM.runFrom, Turing.seamCompTM, Turing.seamCompTM_run_ofCfg}
Let two machines have respective state sets, and suppose the first machine's contract
starts at its entry state and exits exactly at a designated state, while the second
machine's contract starts at its own entry state and exits in a specified set. Given a
map from the first contract's data to the second's and a condition ensuring that the
first endpoint, with its state reset to the second entry state, satisfies the second
precondition, their sequential composition has a contract from the first entry state: it
runs for the first running time plus one transition to switch machines plus the second
running time, and ends at the second endpoint with its state mapped into the second
summand. Its head-position bounds are the coordinatewise minimum and maximum of the two
contracts' bounds.
\end{definition}

\begin{definition}[Binding a configuration's state]
\lean{Turing.Cfg.bindState}
\label{Turing.Cfg.bindState}
\leanok
\uses{Turing.Cfg}
Given a configuration with state set $S$ and a function from $S$ to optional states in
$S'$, form a configuration with state set $S'$ by applying that function to the optional
state component. All other configuration data are unchanged.
\end{definition}

\begin{lemma}[Trajectory through a phase host]
\lean{Turing.Contract.inPhase_traj}
\label{Turing.Contract.inPhase_traj}
\leanok
\uses{Turing.Cfg, Turing.Cfg.bindState, Turing.Cfg.mapState, Turing.Contract, Turing.MultiTapeTM, Turing.MultiTapeTM.runFrom, Turing.phaseHostTM, Turing.phaseHostTM_run}
\difficulty{5}
Let $I$ index phase state sets $\mathrm{Ph}(i)$ and machines $M_i$, let $q_0$ be a state
in their disjoint union, and let $red$ assign to each phase and state either no
successor or a state in that union. Fix a phase $p$, a state $s_0\in\mathrm{Ph}(p)$, a
contract $C$ for $M_p$ with exit set $X$, a datum $i$, and a host configuration $d$
whose state is $(p,s_0)$. Assume that the halted state is not in $X$, every state
outside $X$ reduces to no successor in phase $p$, every state in $X$ reduces to a state
in a different phase, and the configuration obtained from $d$ by replacing its state
with $s_0$ satisfies $C$'s precondition. Through and including $C$'s assigned time, the
host run agrees with the run of $M_p$ from that state-replaced configuration, with
states embedded in phase $p$. One step later, the host run equals the endpoint of $C$
with its state bound through $red$; moreover, that endpoint has a state belonging to
$X$.
\end{lemma}

\begin{definition}[Contract for a phase transition]
\lean{Turing.Contract.inPhase}
\label{Turing.Contract.inPhase}
\leanok
\uses{Turing.Cfg.bindState, Turing.Cfg.mapState, Turing.Contract, Turing.Contract.inPhase_traj, Turing.Contract.mapState_state, Turing.MultiTapeTM, Turing.phaseHostTM}
Let a family of machines have state spaces indexed by a type $I$, and let a reduction
rule optionally send a state in phase $i$ to a state in another phase. Suppose a
contract for the machine in phase $p$ starts at $s_0$ and exits in a set $X$, with no
halted state in $X$; every state outside $X$ is unchanged by the reduction rule, and
every state in $X$ is sent to some different phase. Then the phase-host machine, which
applies that reduction rule when available and otherwise runs the machine for the
current phase, has a contract starting at $(p,s_0)$: its time is the given contract's
time plus one, its endpoint state is obtained by applying the reduction rule to the
original endpoint state, and its exit set consists of those reduced states arising from
states in $X$.
\end{definition}

\begin{lemma}[Recovery after state relabeling]
\lean{Turing.Contract.mapState_const_inl_inl}
\label{Turing.Contract.mapState_const_inl_inl}
\leanok
\uses{Turing.Cfg, Turing.Cfg.mapState}
\difficulty{3}
Let $d$ be a configuration whose state is the left-left injection of some $e\in S_0$.
Relabel every state of $d$ by the constant value $e$, then inject the resulting state
first into $S_0\oplus S_1$ and then into $(S_0\oplus S_1)\oplus S_2$. The resulting
configuration is $d$.
\end{lemma}

\begin{lemma}[Run of a two-branch composition]
\lean{Turing.Contract.branch_traj}
\label{Turing.Contract.branch_traj}
\leanok
\uses{Turing.Cfg, Turing.Cfg.mapState, Turing.Contract, Turing.Contract.mapState_const_inl_inl, Turing.MultiTapeTM, Turing.MultiTapeTM.runFrom, Turing.MultiTapeTM.runFrom_state_ne_none_of_le, Turing.seamCompTM, Turing.seamCompTM_runFrom_left, Turing.seamCompTM₂, Turing.seamCompTM₂_run_ofCfg}
\difficulty{5}
Let $M_0,M_1,M_2$ be machines with state sets $S_0,S_1,S_2$, let $e_1,e_2\in S_0$ be
distinct, and let $C_0$ be a contract for $M_0$ with entry state $entry_0$ and exits
$e_1,e_2$. Given branch entry states $entry_1,entry_2$ and a datum and starting
configuration satisfying the $C_0$ precondition and entry-state condition, the composed
machine's run agrees, through and including the assigned $C_0$ time, with the $M_0$ run
whose states are embedded in the $S_0$ summand; throughout this interval that $M_0$ run
is not halted. If the $C_0$ endpoint has state $e_1$, then after one switching step the
composed run, at every subsequent time, agrees with the run of $M_1$ started from that
endpoint with its state replaced by $entry_1$, with states embedded in the first branch.
If the endpoint has state $e_2$, the corresponding equality holds for $M_2$, started
with the endpoint state replaced by $entry_2$ and states embedded in the second branch.
\end{lemma}

\begin{definition}[Branching composition of contracts]
\lean{Turing.Contract.branch}
\label{Turing.Contract.branch}
\leanok
\uses{Turing.Cfg.mapState, Turing.Contract, Turing.Contract.branch_traj, Turing.Contract.live, Turing.Contract.mapState_state, Turing.MultiTapeTM, Turing.seamCompTM₂}
Let a first machine's contract start at an entry state and exit at one of two distinct
designated states. For each exit state, suppose there is a contract for a corresponding
second machine, together with a map from the first contract's data to that contract's
data; assume the first endpoint, with its state reset to the corresponding second entry
state, satisfies that second contract's precondition whenever it reaches that exit. The
machine that runs the first machine and then switches to the selected second machine has
a contract whose exit set is the union of the two second exit sets mapped into their
respective state summands; its running time is the first running time plus one switching
step and the running time of the selected second contract, and its head-position bounds
are the coordinatewise minimum and maximum of the first bounds and the selected second
bounds.
\end{definition}

\begin{lemma}[Releasing the distinguished start state]
\lean{Turing.release_start}
\label{Turing.release_start}
\leanok
\uses{Turing.Cfg, Turing.Cfg.mapState}
\difficulty{3}
Let $k$ be a natural number, let $\Gamma$ and $S$ be types, choose $\mathrm{anchor}\in
S$, and let $d$ be a configuration on an input word $x$ whose states lie in the disjoint
union of a singleton and $S$. If the state of $d$ is the left summand's unique element,
then first mapping every state to $\mathrm{anchor}$ and then mapping every state to that
left-summand element returns the original configuration $d$.
\end{lemma}

\begin{definition}[Releasing a return contract]
\lean{Turing.ReturnContract.release}
\label{Turing.ReturnContract.release}
\leanok
\uses{Turing.Cfg.mapState, Turing.Contract, Turing.MultiTapeTM, Turing.MultiTapeTM.runFrom_zero, Turing.ReturnContract, Turing.release_start, Turing.seamReleaseTM, Turing.seamReleaseTM_firstReturn, Turing.seamReleaseTM_runFrom_pos}
Given a return contract for a machine that starts and ends at an anchor state, with
positive running time and no return to the anchor at any strictly intermediate positive
time, the release machine has a contract starting in its initial release state and
exiting at the anchor state in the machine-state summand. The released contract inherits
the precondition, running time, endpoint (with its state mapped into that summand), and
head-position bounds from the return contract.
\end{definition}

\begin{definition}[Releasing a round contract]
\lean{Turing.RoundContract.release}
\label{Turing.RoundContract.release}
\leanok
\uses{Turing.Cfg.mapState, Turing.Contract, Turing.MultiTapeTM, Turing.MultiTapeTM.runFrom, Turing.MultiTapeTM.runFrom_zero, Turing.RoundContract, Turing.release_start, Turing.seamReleaseTM, Turing.seamReleaseTM_runFrom_pos}
Given a round contract for a machine that starts at an anchor, has positive running
time, and ends either at the anchor or in an escape set, the release machine has a
contract starting in its initial release state. Its endpoint is the round endpoint with
its state mapped into the machine-state summand, its exit set is the image of the union
of the anchor state and the escape set under that state map, and its precondition,
running time, and head-position bounds come from the round contract.
\end{definition}

\begin{theorem}[Trajectory property for sequential contracts]
\lean{Turing.Contract.Traj.seq}
\label{Turing.Contract.Traj.seq}
\leanok
\uses{Turing.Cfg, Turing.Cfg.mapState, Turing.Contract, Turing.Contract.Traj, Turing.Contract.seq, Turing.Contract.seq_traj, Turing.MultiTapeTM}
\difficulty{5}
Let $C_1$ be a contract for a machine $M_1$ from entry state $e_1$ to the sole exit
state $f_1$, and let $C_2$ be a contract for a machine $M_2$ from entry state $e_2$ to
an exit set $X_2$, on the same input and with the same number of work tapes. Suppose a
map links the data for $C_1$ to the data for $C_2$ and every eligible endpoint of $C_1$,
with its state replaced by $e_2$, satisfies the corresponding $C_2$ precondition. If
$R_1$ and $R_2$ hold throughout the trajectories of $C_1$ and $C_2$, respectively, then
the sequentially composed contract satisfies, at every time, the disjunction of $R_1$ on
the initial phase and $R_2$ on the linked datum and the intermediate endpoint with its
state replaced by $e_2$.
\end{theorem}

\begin{theorem}[Frames under sequential composition]
\lean{Turing.Contract.Frames.seq}
\label{Turing.Contract.Frames.seq}
\leanok
\uses{Turing.Cfg, Turing.Cfg.SameOutside, Turing.Cfg.SameOutside.mapState, Turing.Cfg.SameOutside.trans, Turing.Cfg.mapState, Turing.Contract, Turing.Contract.Frames, Turing.Contract.seq, Turing.Contract.seq_traj, Turing.MultiTapeTM}
\difficulty{4}
Let $M_1$ and $M_2$ be machines with the same input alphabet and number of work tapes,
with state sets $S_1$ and $S_2$, and suppose $S_1$ has decidable equality. Let $C_1$ be
a contract for $M_1$ from entry state $e_1$ to the singleton exit state $e_1'$, and let
$C_2$ be a contract for $M_2$ from entry state $e_2$ to an exit set $X_2$. Let
$\mathrm{link}$ map the data for $C_1$ to data for $C_2$, and suppose that whenever a
starting configuration is in state $e_1$ and satisfies the precondition of $C_1$, the
endpoint of $C_1$ with its state reset to $e_2$ satisfies the corresponding precondition
of $C_2$. If $C_1$ and $C_2$ have the frame property for frame sets $fp_1$ and $fp_2$,
respectively, then their sequentially composed contract has the frame property for the
union of $fp_1$ at the starting configuration and $fp_2$ at the linked data and
state-reset endpoint of $C_1$.
\end{theorem}

\begin{theorem}[Trajectory property within a phase]
\lean{Turing.Contract.Traj.inPhase}
\label{Turing.Contract.Traj.inPhase}
\leanok
\uses{Turing.Cfg, Turing.Cfg.mapState, Turing.Contract, Turing.Contract.Traj, Turing.Contract.inPhase, Turing.Contract.inPhase_traj, Turing.MultiTapeTM}
\difficulty{5}
Let $I$ index a family of machines with state sets $\mathrm{Ph}(i)$, and let a reduction
rule optionally send a state in phase $i$ to a state in one of the phases. Fix a phase
$p$, a contract $C$ for its machine, and a trajectory predicate $R$. Suppose the halted
state is not in the exit set of $C$, every state outside that exit set is left unchanged
by the reduction rule in phase $p$, and every state in the exit set is sent to a state
in a phase different from $p$. If $R$ holds throughout the trajectories of $C$, then it
holds throughout the trajectories of the phase-hosted contract for the configuration
formed by replacing the starting configuration's state with the entry state of $C$, with
the same head-position vector.
\end{theorem}

\begin{theorem}[Trajectory property for branching contracts]
\lean{Turing.Contract.Traj.branch}
\label{Turing.Contract.Traj.branch}
\leanok
\uses{Turing.Cfg, Turing.Cfg.mapState, Turing.Contract, Turing.Contract.Traj, Turing.Contract.branch, Turing.Contract.branch_traj, Turing.Contract.mapState_state, Turing.MultiTapeTM}
\difficulty{5}
Let $M_0,M_1,M_2$ be machines, with distinct states $e_1,e_2$ of $M_0$, and let $C_0$ be
a contract for $M_0$ whose exit set consists of $e_1$ and $e_2$. Let $C_1,C_2$ be
contracts for $M_1,M_2$, respectively, with data maps from $C_0$'s data. Suppose that
whenever the endpoint of an eligible $C_0$ run has state $e_1$ (respectively, $e_2$),
the endpoint with its state replaced by the entry state of $C_1$ (respectively, $C_2$)
satisfies that contract's precondition. If $R_0,R_1,R_2$ hold throughout the
trajectories of their respective contracts, then the branching contract satisfies, at
every time, the disjunction of $R_0$ on the initial phase and the property for the
selected branch: $R_1$ if the $C_0$ endpoint has state $e_1$, and $R_2$ otherwise.
\end{theorem}

\begin{theorem}[Sequential Composition Has Positive Duration]
\lean{Turing.Contract.seq_positive}
\label{Turing.Contract.seq_positive}
\leanok
\uses{Turing.Cfg.mapState, Turing.Contract, Turing.Contract.Positive, Turing.Contract.seq, Turing.MultiTapeTM}
\difficulty{1}
Let $M_1,M_2$ be machines with the same number of work tapes and input alphabet, and
state sets $S_1,S_2$. Let $C_1$ be a contract for $M_1$ from entry state $e_1$ with exit
set $\{\mathrm{some}(f_1)\}$, and let $C_2$ be a contract for $M_2$ from entry state
$e_2$ with exit set $X_2$. Let $\ell$ map the data values of $C_1$ to those of $C_2$,
and suppose that for every admissible $C_1$ start in state $e_1$, the endpoint of $C_1$,
with its state replaced by $e_2$, satisfies the $C_2$ precondition at data value
$\ell(i)$. Then the sequentially composed contract, which runs $C_1$, takes one
transition to switch machines, and then runs $C_2$, has strictly positive assigned
running time for every admissible start in its entry state.
\end{theorem}

\begin{theorem}[Positivity of a phase contract]
\lean{Turing.Contract.inPhase_positive}
\label{Turing.Contract.inPhase_positive}
\leanok
\uses{Turing.Contract, Turing.Contract.Positive, Turing.Contract.inPhase, Turing.MultiTapeTM}
\difficulty{1}
Let $I$ index a family of machines with phase state sets $\mathrm{Ph}(i)$, let $q_0$ be
a state of the phase-host machine, and let $red$ assign to each phase and state either
no successor or a state in the disjoint union of the phase state sets. Fix a phase $p$,
a state $s_0\in\mathrm{Ph}(p)$, and a contract $C$ for the machine in phase $p$. Suppose
the halted state is not in $C$'s exit set; every state outside that exit set reduces to
no successor in phase $p$; and every state in that exit set reduces to a state in a
different phase. If $C$ has strictly positive assigned time for every eligible datum and
starting configuration, then the contract obtained by placing $C$ inside the phase-host
machine has strictly positive assigned time for every eligible datum and starting
configuration.
\end{theorem}

\begin{theorem}[Positivity of a branched contract]
\lean{Turing.Contract.branch_positive}
\label{Turing.Contract.branch_positive}
\leanok
\uses{Turing.Cfg.mapState, Turing.Contract, Turing.Contract.Positive, Turing.Contract.branch, Turing.MultiTapeTM}
\difficulty{1}
Let $M_0,M_1,M_2$ be machines, with $C_0$ a contract for $M_0$ from entry state
$entry_0$ to the two exits $e_1,e_2$, where $e_1\ne e_2$, and let $C_1,C_2$ be contracts
for $M_1,M_2$. Let $l_1,l_2$ map data for $C_0$ to data for $C_1,C_2$, respectively.
Suppose that whenever an admissible $C_0$ endpoint has state $e_1$ (respectively,
$e_2$), the endpoint with its state replaced by the entry state of $M_1$ (respectively,
$M_2$) satisfies the corresponding branch contract's precondition. Then the branched
contract has strictly positive assigned time for every datum and starting configuration
satisfying its entry-state and precondition conditions.
\end{theorem}

\begin{theorem}[Positive Duration of the Released Contract]
\lean{Turing.ReturnContract.release_positive}
\label{Turing.ReturnContract.release_positive}
\leanok
\uses{Turing.Cfg.mapState, Turing.Contract.Positive, Turing.MultiTapeTM, Turing.ReturnContract, Turing.ReturnContract.release}
\difficulty{2}
Let $M$ be a multi-tape Turing machine, $a$ an anchor state, and $C$ a return contract
for $M$ at $a$. For every data value $i$ and configuration $d$ whose state is the
release machine's initial state and which satisfies the released contract's
precondition, the released contract's prescribed running time is strictly positive.
\end{theorem}

\begin{theorem}[Positive Duration of Released Rounds]
\lean{Turing.RoundContract.release_positive}
\label{Turing.RoundContract.release_positive}
\leanok
\uses{Turing.Cfg.mapState, Turing.Contract.Positive, Turing.MultiTapeTM, Turing.RoundContract, Turing.RoundContract.release}
\difficulty{2}
Let $M$ be a multi-tape Turing machine with anchor state $a$, escape set $E$, and round
contract $R$. For every data value $i$ and configuration $d$ whose state is the release
machine's initial state and which satisfies the released contract's precondition, the
released contract's prescribed running time is strictly positive.
\end{theorem}
